{
\psset{linearc=0}
{\Large Як читати цю діаграму}
\begin{itemize}

\item Ця діаграма організована в октавах (подвоєння/половинення частоти), починаючи з 1Hz і зростаючи (2, 4, 8 тощо) та зменшуючись (1/2, 1/4 тощо). Октава --- це природний спосіб представлення частоти.

\item Частота зростає у напрямку вгору і переходить справа наліво.

%\item This chart was also designed so that the spectrum section could be cut our and wrapped around a standard 3 inch diameter by 36 inches long mailing tube to represent a continuous frequency from the bottom to the top and beyond.

\item Ця діаграма не має обмежень з жодного боку, проте через обмежений простір тут показано лише ``відомі'' елементи. Частота 0Hz є найнижчою можливою частотою, але метод зображення октав, використаний тут, не дозволяє досягти 0Hz, лише наближатися до неї. Також, за визначенням частоти (цикли за секунду), від'ємної частоти не існує.

\item Позначення значень: \psframebox[framesep=1pt,fillstyle=solid,fillcolor=Black]{\textcolor{FColor}{Частота}} в герцах,%
\psframebox[framesep=1pt,fillstyle=solid,fillcolor=Black]{\textcolor{WColor}{Довжина хвилі}} в метрах,%
\psframebox[framesep=1pt,fillstyle=solid,fillcolor=Black]{\textcolor{EColor}{Енергія}} в електронвольтах.

\end{itemize}
}
